% !TeX root = ../../../main.tex
\section{Methods}

\subsection{Patient Recruitment}

All procedures performed in this study were approved by the Institutional Review Board at the University of North Carolina at Chapel Hill (UNC), and written informed consent was obtained from all participants prior to enrollment (ClinicalTrial.gov No. NCT03079882). Patients over 18 years of age who received renal transplants at the UNC Medical Center and underwent follow-up care through the UNC healthcare network (UNC Health) were recruited within three months post-transplantation for this prospective, longitudinal case-control study. Patients whose renal allografts were not visible at a maximum imaging depth of 40.0 mm based on B-mode imaging using the equipment, settings, and protocols described in Section~\ref{sec:transplantImageAcq}, as well as those with contraindications to ultrasound imaging (e.g. open wounds, infections, staples at the imaging site) or diagnoses of graft rejection prior to recruitment were excluded from participation. Patients were imaged during routine onsite clinical follow-up visits at the kidney transplant clinic at UNC Health. This regimen continued for up to patients' followup visit marking three years post-transplantation, until clinical indication for biopsy, or until loss to follow-up.\footnote{
    This study was originally designed to allow for images of each patient to be acquired at scheduled clinical visits every 3 months up to three years post-transplantation such that VisR parameters could be regularly acquired and correlated with changes in clinical status over time. However, this imaging schedule was significantly disrupted by logistical challenges noted in Section~\ref{sec:transplantDiscussionLimitations}. 
} All patient care decisions were made independently from study participation, and imaging was performed separately from clinical workflows. At each clinica
l visit including the one where patients were recruited, demographic information (e.g. BMI) and clinical laboratory results (e.g. serum creatinine, urinary protein-creatinine ratio) were collected from patients' electronic health records as part of standard of care.

\subsection{Image Acquisition and Processing}
\label{sec:transplantImageAcq}

At each scheduled clinical visit, recruited patients underwent VisR imaging of their renal allografts as performed by an experienced sonographer who was blinded to patients' clinical histories and laboratory results. Patients were imaged in the supine position on a bed with a 15\textdegree{} incline and asked to remain motionless during imaging. Based on the protocols described in \cite{hossain_Acoustic_2017,hossain_Evaluating_2018}, and \cite{hossain_Variability_2019}, VisR imaging was performed using a VF7-3 linear array and Acuson Antares clinical-grade scanner (Siemens Healthineers, Issaquah, WA, USA) modified with a research interface to allow for custom beam sequencing and raw RF data acquisition.

Each patient's allograft was imaged three times in each of two orthogonal imaging planes. First, the long direction of the linear transducer array was aligned with the long axis of the kidney, and the superior pole of the kidney was identified and imaged such that (1) the cortex and medulla were visible and (2) large vessels were avoided. This transducer orientation, denoted as the "0 degree" or "along" orientation, was used to acquire the first set of three VisR acquisitions. In the event where the superior pole was not accessible due to patient anatomy or other factors, the transducer was placed such that the long axis of the transducer was still aligned with the long axis of the kidney but placed to image the inferior pole. When the transducer position was determined, an appropriate focal depth was selected such that nephrons in the renal cortex were expected to be parallel to the transducer surface based on B-mode imaging. With the transducer orientation and position as well as focal depth fixed, three VisR acquisitions were taken in this orientation. If the medulla could be better visualized by translating the probe laterally (i.e. along the long axis of the kidney, closer to the renal pelvis), three additional images were taken at that location. Then, at the original position closer to the chosen pole of the kidney, the transducer was rotated 90 degrees such that the long axis of the transducer was perpendicular to the long axis of the kidney (i.e. to the "90 degree" or "across" orientation). Three additional VisR acquisitions were taken in this orientation, as well, leading to a maximum total of nine VisR acquisitions per patient per clinical visit using one focal depth.

For each VisR acquisition, two Hamming-apodized, 300-cycle (71 \unit{\micro\second}) ARF excitations with a center frequency of 4.21 MHz and F/1.5 focal configuration were applied. Induced displacements were tracked on-axis to the ARFI pushes using 2-cycle pulses with a Hamming apodization at 6.15 MHz, repeating at a PRF of 11.5 kHz. Two tracking pulses were emitted prior to the first ARFI push, followed by 6 tracking pulses between pushes and 66 pulses after the second push. This ensemble was repeated across a 20.0 mm lateral field of view. Finally, co-localized B-mode images were acquired using conventional focused imaging with a center frequency of 6.15 MHz and an F/1.5 focal configuration with the maximum 40.0 mm lateral field of view allowed by the scanner for ROI placement. Hardware-based F/0.75 dynamic receive focusing was automatically applied during tracking and B-mode imaging.

\subsection{Evaluation of VisR Parameters and Semi-Quantitative Metrics}

The resulting raw RF data, acquired at a 40 MHz sampling rate, was motion-tracked into displacement profiles~\cite{pinton_Rapid_2006} that denote one-dimensional axial displacements of tissue over time at each lateral location and depth. The resulting displacement profiles were then motion-filtered~\cite{giannantonio_Comparison_2011} using a quadratic motion filter to mitigate the effects of exogenous motion (e.g. breathing, peristalsis). First, for each lateral location and depth, the peak displacement over time following by the second ARFI push (PD) and the logarithm of the variance of acceleration (VoA) of the final 20 samples of displacement estimates~\cite{torres_ARFI_2017,torres_Delineation_2019} were calculated. Then, each displacement profile was fit to the MSD model using a Nelder-Mead nonlinear optimizer~\cite{selzo_Quantitative_2016} implemented in Matlab R2019a (MathWorks, Natick, MA, USA). From the fitted MSD parameters, RE and RV were calculated at each lateral location and depth using Eqs.~\eqref{eqRe} and \eqref{eqRv}, respectively. 

For each acquisition, a trained reviewer who was blinded to patients' clinical histories and laboratory results reviewed the co-localized B-mode, PD, and log(VoA) images. Given the option to toggle between overlays of PD and log(VoA) images on top of the B-mode for guidance, the reviewer placed 2.0 mm square regions of interest (ROIs) in the outer cortex, inner cortex, outer medulla, and inner medulla of the renal parenchyma. ROIs were only placed within a 4.0 mm axial range centered at the ARF push focal depth to mitigate depth-dependent variations in ARF amplitude. When anatomical variations, imaging artifacts, or limited visibility precluded the placement of ROIs in any of the four anatomical regions, ROIs were not placed in those regions. Given ROI positions for each acquisition, median values of PD, RE, and RV within each ROI were calculated. Note that log(VoA) images were only used to guide ROI placement, and were not used for quantitative measurements.

For each patient, clinical visit, transducer orientation, and anatomical region, up to three PD, RE, and RV measurements were taken; where multiple measurements existed for the same anatomical region in the same orientation (e.g. when an imaging plane intended for the cortex captured a second ROI for the outer medulla), the 25th, 50th, and 75th percentiles of those measurements were calculated using the Matlab \texttt{prctile(..., [25, 50, 75])} function to represent the distribution of measurements. Using PD, RE, and RV measurements from both the 0 degree and 90 degree orientations, DoA values were calculated using Eq.~\eqref{eqDoA} for all possible combinations of ROIs where measurements existed in both orientations for the same anatomical region, and the aforementioned percentile binning approach was applied to calculate three representative samples for the distribution of DoA values. The same approach was used to pair 0 degree measurements from two different anatomical regions (e.g. outer cortex and inner cortex) to calculate RR values using Eq.~\eqref{eqRr}.

\subsection{Histological Comparisons and Statistical Analyses}

As UNC Health does not perform protocol biopsies for renal allograft monitoring, graft health outcomes were assessed based on biopsies performed in two situations as part of standard of care: (1) due to zero-hour biopsies performed at the time of transplantation at the discretion of the surgical team, and (2) due to clinical suspicion of graft dysfunction such as increases in serum creatinine, detection of donor-specific antibodies (DSA), and delayed graft function. Biopsy slides were acquired, processed, and scanned as part of routine clinical care by board-certified renal pathologists at UNC Health who were blinded to VisR imaging results. The slides were accessed retrospectively and evaluated according to the 2022 Banff classification criteria~\cite{naesens_Banff_2024} by board-certified nephropathologists who were blinded to VisR imaging results.\footnote{
    Data acquisition for this study began in late 2016, when the widely accepted procedure for renal allograft pathology evaluation was the 2015 Banff criteria~\cite{loupy_Banff_2017}. However, this study continued through early 2021, leading to a situation where biopsies may have been evaluated according to the 2015 standards or updates from the biannual meeting in 2017~\cite{haas_Banff_2018} or 2019~\cite{loupy_Banff_2020} depending on when the biopsy was performed. This posed a challenge for consistent comparisons across the study period, as the definition of certain terminologies such as "suspicious for [AMR or TCMR]", "chronic active TCMR", and "borderline changes" were introduced, modified, or removed over time. To standardize biopsy evaluations across the study period, all biopsy slides were re-evaluated according to the 2022 criteria~\cite{naesens_Banff_2024}, the most recent iteration at the time of analysis. It should be noted that the 2024 update~\cite{naesens_Banff_2025}, which reaffirmed the 2022 criteria, was undergoing preliminary review by the nephrology community at the time of analysis, and had not been finalized or submitted to any peer-reviewed publications at the time of writing. Thus, the 2022 criteria were used for biopsy re-evaluations.
} Specifically, biopsy slides were graded on the applicability and severity of interstitial fibrosis and tubular atrophy (IFTA) diagnostic category, lesion scores for interstitial inflammation (\emph{i} score) and inflammation co-present in areas with IFTA (\emph{i-IFTA} score), and for the presence of antibody-mediated rejection (AMR) or T-cell mediated rejection (TCMR). If a patient underwent multiple biopsies during their study participation, all biopsy results from their participation were recorded for analysis. Each VisR imaging date was associated with the most recent indication biopsy. If a patient did not undergo a qualifying indication biopsy within that window for their first VisR imaging date but did undergo a zero-hour biopsy, the results of the zero-hour biopsy were associated with that VisR imaging date.

Statistical analyses were performed using Matlab R2019a. VisR measurements (DoA and RR from PD, RE, and RV) were grouped by time of imaging relative to their final biopsy/imaging date, patient demographics, biopsy findings (\emph{i} score, IFTA grade, \emph{i-IFTA} score), and clinical diagnoses. For each anatomical region (DoA) or pair thereof (RR), differences in VisR measurements  across groups were assessed using the Kruskal-Wallis test with Bonferroni post-hoc correction with a threshold of $\alpha$ = 0.05. Correlations among demographic and clinical factors were assessed using Spearman's rank correlation coefficient using the \texttt{corr(..., 'Type', 'Spearman', 'Rows', 'complete')} function in Matlab R2019a. The significance of correlations were assessed using the corresponding $p$-values outputted by the aforementioned function with a threshold of $\alpha$ = 0.05.

The ability to predict clinical variables from VisR measurements was assessed using two approaches. For continuous variables and ordinal variables with more than two classes (e.g. IFTA severity, which has four levels from 0 to 3), regression models were trained to directly predict the same variables. Additionally, for all variables, binary classifiers were trained to identify thresholds of clinical variables. For regressors that were made to predict continuous clinical variables, linear regression models were created using a least-squares approach, and binary classifiers were created based on commonly accepted clinical thresholds using logistic regression~\cite{pagana_Creatinine_2021,bouchard_Estimating_2021,haider_Proteinuria_2025}. For supported categorical clinical variables, mulinomial probit regression models were created using the \texttt{mnrfit} function, while multiclass classifiers were developed using the \texttt{fitcecoc} function with a one-versus-one coding design. For all other (i.e. binary clinical variables), weighted logistic regression was applied. All models were developed using a $k$-fold nested cross-validation scheme according to \cite{varma_Bias_2006}, where $k = \min \left( 10, k_{\textrm{min}} \right)$ and $k_{\textrm{min}}$ is the minimum number of samples in any class of the clinical variable being predicted. The performance of continuous response variable-based regressors were assessed by applying a one-sample $t$-test to evaluate (1) whether Spearman's rank correlation coefficients following a Fisher's $z$ transformation were zero, and (2) whether the RMSE of the model was not significantly lower than that of a naive regressor that predicts the mean value of the clinical variable for all samples associated with a biopsy and image. For categorical response variable-based regressors, the latter null hypothesis instead evaluated whether the cross-entropy loss of the model was not significantly lower than that of a random classifier. For classification models, the hypothesis tests evaluated whether (1) the area under the receiver operating curve (AUC), (2) sensitivity, and (3) specificity were not significantly greater than those of random guessing following a Fisher's $z$ transform. The full list of variables considered, as well as the model types and performance metrics used for evaluation, are listed in Table~\ref{tbl:clinicalVariables}.

% !TeX root = ../../../main.tex

\begin{table}[tbp]
    \centering
    \caption{Model Definitions and Assessment Approaches for Regression and Classification Models Predicting Clinical Variables from VisR-Derived Metrics}
    \begin{tabular}{llrll}
        \toprule
        \multirow{2}{*}{Parameter} & 
        \multicolumn{2}{c}{Regressor} & 
        \multicolumn{2}{c}{Classifier}\\
         & Metric & Baseline & Model Type & Categories\\
        \midrule
        Serum Creatinine (mg/dl) &
        RMSE & 0.447 &
        Binary & $<$ 1.6, $\geq$ 1.6\\
        UPCR (g/g) &
        RMSE & 0.589 &
        Binary & $<$ 0.862, $\geq$ 0.150\\
        DSA Presence &
        - & - &
        Binary & Absent, Present\\
        Indication Biopsy &
        - & - &
        Binary & No, Yes\\
        \emph{i} Score &
        - & - &
        Binary & 0, 1--3\\
        IFTA Grade &
        Cross-Entropy & 1.39 &
        Multinomial & 0, 1, 2, 3\\
        \emph{i-IFTA} Score &
        Cross-Entropy & 1.10 &
        Multinomial & 0, 1, 2\\
        Rejection &
        - & - &
        Binary & No, Yes\\
        \bottomrule
    \end{tabular}
    \par\vspace{0.5em}
    \noindent{}RMSE: root mean square error. Baseline metrics for RMSE are calculated as the root mean square of the difference between individual values and their medians. Baseline metrics for cross-entropy losses are based on naive models that randomly predict each category with equal probability. For example, if a clinical parameter has three categories, the baseline is $-ln \left( 1 / 3 \right) \approx 1.10$.
    \label{tbl:clinicalVariables}
\end{table}

